\section{Language models as the arithmetic unit}\label{sec:llm}

In Zuse's division of roles, the language model is the arithmetic unit. It formulates statements, writes requests and reads receipts.
What counts as a result is decided by the record and the executor. Two mechanisms keep the roles apart.

\subsection{Read-only access through MCP}

Language models reach the record through the Model Context Protocol~\cite{mcp}, a standard interface through which a model calls
tools. QAM provides servers with graded rights. The read-only server lets a model search entries, list and retrieve receipts (with
their premise trace or as a summary), select receipts by context, and list, read, compare and audit batch reports. It is meant to be run, as documented, in a
container without network access, with a read-only file system and the database mounted read-only. Further servers can rebuild the index and execute
single requests or batches; they add receipts and offer no operation that changes an existing one. All servers run only fixed modules
under a time limit. None of them offers a tool that raises a status, rewrites a record or deletes a receipt.

A model asked about a stored result is expected to retrieve the receipt, verify it by its hash and state its open guards, rather than
rely on an earlier conversation. The receipt summary is designed for this: it gives the verified hash, the instruction, the outcome,
errors, obligations and the declared guards with their status in a few lines.

\subsection{Hallucination as a status error}

QAM does not detect hallucinations as errors of content; it could not. It treats them as \emph{status errors}: as attempts to state
something more strongly than its premises allow, or to present an input as a result. This kind of error is mechanically detectable,
because it shows in the structure of the record rather than in the meaning of the text. A language model can omit a condition from the
prose. It cannot remove it from the record without changing hashes, and it cannot set the status of a conclusion above that of its
premises without the compiler (in its default mode) or the Lean checker rejecting the record. A calculated value, finally, is backed
by evidence only if the calculation was submitted as a request and left a receipt; a value merely declared \st{computed} in the record
has no receipt that could be retrieved for it.

This does not make a model reliable. It changes where its unreliability can do damage: from the silent transitions of
Section~\ref{sec:intro}, which are hard to see, to explicit entries and receipts, which can be checked.
